\section{Materials and methods}\label{materials-and-methods}

\subsection{Dataset and experimental setup}\label{dataset-and-experimental-setup}

The experiments were conducted on a Blohm Orbit 36 precision surface-grinding machine tool. The grinding process used a vitrified-bond cubic boron nitride (cBN) grinding wheel (specification B126-V240, diameter 400~mm). The workpiece material was hardened AISI 52100 bearing steel (approximate hardness 60~HRC). A water-soluble 5\% emulsion coolant was applied at a nozzle pressure of 2.5~bar and a flow rate of 40~L/min. During each grinding pass, three modalities were acquired simultaneously:

\begin{itemize}
\tightlist
\item
  \textbf{Acoustic emission (AE):} a two-channel AE acquisition sampled at
  $4.0$~MHz. The raw AE files contain two physical acquisition columns:
  column~1 is labelled narrowband AE and column~2 broadband AE in
  \texttt{data/AE meta.txt}; these map respectively to AE planes~1 and~2 in
  the cached spectrograms. Both physical channels are retained in every
  AE-dB and AE-dB-z configuration. The archive does not preserve sensor
  models, analogue filter bands/orders, gains, or calibration records, so
  the labels establish channel provenance but not a calibrated transfer
  function. The two archived channels provide empirical high-frequency
  descriptors of contact-related activity. Interpretation of broadband
  content (0.1--2.0~MHz) remains sensor- and mounting-dependent.
\item
  \textbf{Vibration:} a tri-axial accelerometer sampled at $51.2$~kHz.
  All three axes are retained to capture the full structural response; the Z-axis is shown in the illustrative
  figures because it is perpendicular to the grinding direction and most
  sensitive to wheel--workpiece normal dynamics.
\item
  \textbf{Process parameters:} wheel speed $v_s$ (m/s), workpiece rotational speed $v_w$ (r/min), and grinding depth $a_p$ (\(\upmu\)m). The workpiece was clamped on a rotating index spindle (cylindrical configuration) with an outer diameter $d_w = 214$~mm, converting the rotational speed to a linear surface speed of $v_w \in \{0.224, 0.448, 0.672, 0.896\}$~m/s for kinematic calculations. The 16 grinding conditions correspond to a fractional factorial design crossing four levels of wheel speed $v_s \in \{25, 30, 35, 40\}$~m/s, four levels of workpiece speed $v_w \in \{20, 40, 60, 80\}$~r/min, and four levels of grinding depth $a_p \in \{20, 30, 40, 50\}$~\(\upmu\)m.
\item
  \textbf{Target:} post-process surface roughness $R_a$ measured with a
  stylus profilometer and reported in \(\upmu\)m.
\end{itemize}

Measurement and acquisition details are as follows. The stylus
profilometer measurements (stylus tip radius of 2~\(\upmu\)m, expanded uncertainty of 0.025~\(\upmu\)m, with a new stylus to minimize wear effects) used a 0.8~mm Gaussian cutoff and a 4~mm
evaluation length; three traces were taken on each ground surface and
averaged to obtain the reported $R_a$. The profilometer was calibrated
before each measurement session with a certified roughness standard; the
manufacturer-quoted repeatability is 0.01~\(\upmu\)m. Surface roughness values
in this study range from 0.066 to 0.381~\(\upmu\)m, so the best reported MAE
(0.020~\(\upmu\)m) is only about twice the manufacturer-quoted repeatability.
Only the final averaged $R_a$ value per pass is available; the
individual trace profiles and trace-level $R_a$ values were not archived,
so an experimental repeatability or gauge R\&R cannot be computed here.
The reported MAE combines model error and measurement uncertainty;
because trace-level repeatability is unavailable, the latent error
relative to latent roughness cannot be separated from
measurement uncertainty. This limitation is especially relevant because rank
differences among the top RF variants are smaller than 0.001~\(\upmu\)m,
far below the manufacturer-quoted repeatability. Future work should
retain trace-level $R_a$ values and raw profiles for gauge R\&R.
Consequently, improvements smaller than approximately 0.005--0.010~\(\upmu\)m
should be interpreted cautiously because only manufacturer-quoted
repeatability, not experiment-specific repeatability, is available, and
differences among the top-ranked random-forest variants below this range
should not be treated as practically meaningful. The
ranking among the top five RF variants is therefore exploratory.
Historical campaign notes describe a nominal AE response band of
100~kHz--1~MHz and mounting on the workpiece fixture with acoustic
couplant; the underlying sensor specification is not archived. The
tri-axial accelerometer was mounted near the grinding contact zone.
Campaign documentation indicates analogue anti-aliasing before digitisation,
but the filter cut-offs and orders were not archived. Non-archived no-cut recordings were inspected informally but
are not used as reproducible evidence.
The archived campaign metadata do not retain sensor model numbers,
calibration sensitivities, preamplifier gains, analogue filter cut-offs or
orders, or mounting torque/attachment records. These missing sensor-chain
details limit transfer-function interpretation. The frozen archive will
explicitly record their absence; future acquisition protocols must retain
them.

The 16 conditions were run in a single randomised order; the wheel was
dressed before the start of the campaign and coolant flow was held
constant. We therefore cannot separate condition effects from potential
wheel-wear or thermal-drift trends across the run order. A run-order plot
of measured $R_a$ is provided as Supplementary
Figure~\ref{fig:supp-runorder} and should be considered when
interpreting Condition~7 as an outlier. Supplementary Figure~\ref{fig:supp-runorder-diagnostics} overlays best-RF
residuals, AE broadband RMS, and tri-axial vibration RMS on the run-order
plot. Condition~7 does not coincide with a monotonic run-order trend in
sensor RMS, so a simple wheel-wear or thermal-drift explanation is not
supported by the available diagnostics; in other words, the available
run-order diagnostics do not show a simple monotonic drift explanation
for Condition~7. They do not, however, exclude non-monotonic wheel-state
effects such as local wheel loading, dressing variations, or transient
coolant/temperature events.

The campaign nominally contained 20 passes per condition; after
excluding Condition~1, sample~1, Condition~1 contributed 19 valid
model-evaluation samples and the remaining conditions contributed 20.
Each condition corresponds to a unique combination of the three process
parameters: $v_s \in \{25,30,35,40\}$~m/s, $v_w \in
\{20,40,60,80\}$~r/min, and $a_p \in \{20,30,40,50\}$~\(\upmu\)m.
The 16-run design is a fractional factorial design inspired by
L16-style coverage of the factor levels; it is not the full
$4 \times 4 \times 4 = 64$ factorial. The design provides balanced coverage for wheel speed against the other
two factors, but the workpiece-speed/depth-of-cut pairing is not fully
balanced (Supplementary Table~\ref{tab:supp-balance}). The exact
factor-level combinations are listed in Table~\ref{tbl:condition-design}.
Because the design is fractional, two- and three-way interactions among
wheel speed, workpiece speed, and grinding depth are not independently
estimable and may be aliased or under-sampled, which limits
interaction-effect generalisation.

\begin{table}[htbp]
\centering
\caption{Fractional factorial experimental design. The 16 grinding
conditions are unique combinations of wheel speed $v_s$, workpiece speed
$v_w$, and grinding depth $a_p$ selected from the full
$4 \times 4 \times 4$ factorial. Each condition was nominally replicated
20 times; Condition~1 has 19 valid model-evaluation passes after
exclusion of one corrupt/missing AE/vibration file.}
\label{tbl:condition-design}
\begin{tabular}{@{}ccccc@{}}
\toprule
Condition & $v_s$ (m/s) & $v_w$ (r/min) & $a_p$ (\(\upmu\)m) & Passes \\
\midrule
1  & 25 & 20 & 20 & 19 \\
2  & 25 & 40 & 30 & 20 \\
3  & 25 & 60 & 40 & 20 \\
4  & 25 & 80 & 50 & 20 \\
5  & 30 & 40 & 40 & 20 \\
6  & 30 & 60 & 50 & 20 \\
7  & 30 & 80 & 20 & 20 \\
8  & 30 & 20 & 30 & 20 \\
9  & 35 & 60 & 50 & 20 \\
10 & 35 & 80 & 20 & 20 \\
11 & 35 & 20 & 30 & 20 \\
12 & 35 & 40 & 40 & 20 \\
13 & 40 & 80 & 30 & 20 \\
14 & 40 & 20 & 40 & 20 \\
15 & 40 & 40 & 50 & 20 \\
16 & 40 & 60 & 20 & 20 \\
\bottomrule
\end{tabular}
\end{table}

Consequently, samples within the same condition are more similar to each
other than to samples from different conditions. To evaluate
generalisation to unseen process states, all models are assessed with
leave-one-condition-out (LOGO) cross-validation, detailed in
Section~\ref{validation-protocol}. Essential code, experiment definitions,
and small result manifests are available in the PhysicalGrindingFusion
repository. Large raw recordings and intermediate caches are available for
research use on request rather than mirrored in Git.

\subsection{Signal representations}\label{signal-representations}

Rather than feeding raw waveforms directly to the models, we convert AE
and vibration signals into frequency-domain descriptors covering vibration
bands below 25~kHz and AE bands up to the 2~MHz Nyquist limit from the
4~MHz sampling rate. All
representations are pre-computed and cached so that every model is
trained and evaluated on identical inputs. The overall preprocessing
pipeline is shown in Figure~\ref{fig:preprocessing-pipeline}.

\begin{figure}[htbp]
\centering
\resizebox{\linewidth}{!}{%
\begin{tikzpicture}[
  node distance=1.2cm and 1.4cm,
  box/.style={pub/representation, minimum width=2.4cm, minimum height=0.9cm},
  param/.style={pub/base, dashed, fill=pubSlateFill, minimum width=2.2cm, minimum height=0.7cm},
  arrow/.style={pub/arrow},
  splitarrow/.style={pub/flow}
]
\node[box] (raw) {Raw\\waveform};
\node[box, right=of raw] (stft) {Hann-window\\STFT};
\node[box, right=of stft] (lin) {STFT\\magnitude\\spectrum};
\node[box, above right=0.8cm and 1.2cm of lin] (db) {Mean dB\\spectrogram\\(\texttt{ae\_spec},\\\texttt{vib\_spec})};
\node[box, right=2.2cm of db] (z) {Optional per-sample\\z-standardisation\\(\texttt{ae\_logspec},\\\texttt{vib\_logspec})};
\node[box, below right=0.8cm and 1.2cm of db] (mel) {Inverse dB $\rightarrow$ power\\mel bank $\rightarrow$ dB\\(\texttt{ae\_mel},\\\texttt{vib\_mel})};
\node[param, above=0.4cm of stft] (stftp) {$N_{\rm FFT}$, hop};
\node[param, below=0.4cm of mel] (melp) {$M=64$ mel filters};
\draw[arrow] (raw) -- (stft);
\draw[arrow] (stft) -- (lin);
\draw[splitarrow] (lin) -- (db);
\draw[splitarrow] (db) -- (z);
\draw[splitarrow] (db) -- (mel);
\draw[arrow, dashed] (stftp) -- (stft);
\draw[arrow, dashed] (melp) -- (mel);
\end{tikzpicture}%
}
\caption{Preprocessing pipeline for AE and vibration signals. Raw
waveforms are transformed by a Hann-window STFT. The magnitude branch is
converted to dB and averaged over local maps; an optional per-sample
z-standardisation creates the historical \texttt{*\_logspec} caches. The
implemented log-mel branch inverse-converts the pass-mean dB cache to power,
applies the mel filter bank, and converts the result back to dB. All
representations are cached before model training. Because local maps are
averaged in dB before inverse conversion, this branch mel-filters a
geometric-mean power quantity rather than the arithmetic mean of local power
or log-mel maps.}
\label{fig:preprocessing-pipeline}
\end{figure}

\subsubsection{Time--frequency representations}

Short-time Fourier transforms are computed with a Hann window. For the
AE channel, the fast Fourier transform length is $N_{\text{FFT}}^{\rm
AE}=598$ samples and the hop length is $H^{\rm AE}=426$ samples,
yielding 300 frequency bins and 47 time frames per local AE map. For vibration,
$N_{\text{FFT}}^{\rm vib}=512$ and $H^{\rm vib}=426$, yielding 257
frequency bins and 13 time frames per local vibration map. The STFT complex coefficients are converted to magnitude and then to a decibel scale (clipped at $-80$~dB), producing local dB spectrogram maps. The AE cache has two planes labelled ``narrowband'' and ``broadband'' in the archived intermediate files; the available cache-generation code does not preserve the analogue/digital filter provenance needed to calibrate the channels' transfer paths. Raw-file metadata identify plane~1 as physical narrowband AE and plane~2 as physical broadband AE; this establishes channel provenance but not filter calibration. Both planes are included in every AE-dB and AE-dB-z configuration. The vibration arrays retain the three accelerometer axes. The effective frequency resolutions are
approximately $6.69$~kHz/bin for AE and
$100$~Hz/bin for vibration. Because the sampling rates differ by two
orders of magnitude, the stated STFT parameters yield different
effective physical window durations (transform lengths): approximately 150~\(\upmu\)s for AE and
10~ms for vibration, corresponding to total segment durations of approximately 5~ms and 
110~ms, respectively. The cached per-pass descriptor is not a single
middle-pass window. The per-pass intermediate \texttt{*\_spec.npz} files
contain \(N_p=2{,}910\) local maps for each modality and pass; the final
dataset-level \texttt{mean\_specs.npz} cache stores only their arithmetic
pass means with shapes \((319,2,300,47)\) and \((319,3,257,13)\). The retained intermediate cache does not
record the upstream local-map hop, overlap, terminal-window policy, or an
AE--vibration timestamp correspondence; consequently, ``aligned'' is not
claimed at event level. If the reported approximately 14.6~s pass duration
and uniform placement are assumed, 2,910 maps imply a nominal index spacing
of about 5.0~ms. This would imply negligible overlap for a 5~ms AE segment
and approximately 95\% overlap for a 110~ms vibration segment, but these are
arithmetic implications rather than recovered acquisition settings. The
common map index supports corresponding pass-position thirds only. For each pass \(p\), channel \(c\), frequency bin
\(f\), and STFT time frame \(t\), \texttt{compute\_mean\_spectrogram()}
stores
\begin{equation}
\bar S_{p,c,f,t}=\frac{1}{N_p}\sum_{j=1}^{N_p}S_{p,j,c,f,t}.
\label{eq:local-map-mean}
\end{equation}
This is a simple arithmetic mean over the local-map axis, with no temporal
weighting, selection, or learned aggregation. Thus the reported models use
a pass-level mean of local 5~ms AE and 110~ms vibration descriptors. This
averaging suppresses temporal ordering and transient localisation. A
retrospective early/middle/late-third and ordered three-window-concatenation
sensitivity analysis is reported in Supplementary
Table~\ref{tab:supp-pass-position}; fixed-distance windows and streaming
dispersion features remain future extensions. Both representations come
from the same grinding pass, but the fused inputs concatenate pass-level
descriptors with unequal local time extents. They are therefore not an
event-synchronised fusion. For tabular models, the maps are flattened and
concatenated, producing 38,223 features for fused full-resolution dB or dB-z
(\(2\times300\times47 + 3\times257\times13\)) and 8,512 features for fused
64-bin log-mel (\(2\times64\times47 + 3\times64\times13\)). Dimensions of the main
time--frequency representations are summarised in
Table~\ref{tbl:representation-dims}.

\begin{table}[htbp]
\centering
\caption{Dimensions of the principal time--frequency representations.
$N_{\rm FFT}$ is the transform length, $H$ the hop length, $C$ the number
of retained input planes or axes, $F$ the number of frequency bins, and
$T$ the number of time frames. For AE, $C=2$ denotes the physical narrowband and broadband AE channels; analogue filter specifications and gains are unavailable; for vibration, $C=3$ denotes the three accelerometer axes. The \texttt{*\_logspec} caches are per-sample z-standardised versions of the full-resolution dB maps, whereas the \texttt{*\_mel} caches are 64-bin log-mel maps.}
\label{tbl:representation-dims}
\begin{tabular}{@{}lccccc@{}}
\toprule
Representation & $N_{\rm FFT}$ & $H$ & $C$ & $F$ & $T$ \\
\midrule
\texttt{ae\_spec} & 598 & 426 & 2 & 300 & 47 \\
\texttt{vib\_spec} & 512 & 426 & 3 & 257 & 13 \\
\texttt{ae\_logspec} & 598 & 426 & 2 & 300 & 47 \\
\texttt{vib\_logspec} & 512 & 426 & 3 & 257 & 13 \\
\texttt{ae\_mel} & -- & -- & 2 & 64 & 47 \\
\texttt{vib\_mel} & -- & -- & 3 & 64 & 13 \\
\bottomrule
\end{tabular}
\end{table}

The historical \texttt{ae\_logspec} and \texttt{vib\_logspec} names denote
full-resolution dB spectrograms after per-sample, per-plane z-standardisation;
they do not apply an additional logarithm. They retain the 300 AE and
257 vibration frequency bins shown in Table~\ref{tbl:representation-dims}
and are standardised per sample and per channel. No subsequent fold-wise
spectral scaler is used in the RF experiments reported here. The 64-bin
log-mel descriptors are stored separately as \texttt{ae\_mel} and
\texttt{vib\_mel}. They are generated from the pass-mean dB cache by inverse
conversion to linear power, followed by multiplication by $M=64$ mel-spaced
filters covering the full Nyquist band. If $D$ is the pass-mean dB cache, the
implemented output is $10\log_{10}[\varepsilon + M(10^{D/10})]$ with a fixed global floor
$\varepsilon = 10^{-10}$ to avoid numerical underflow.
Because \(D\) is averaged over local maps in dB before this inverse
conversion, \(10^{D/10}\) is a geometric-mean power quantity. The implemented
cache is therefore distinct from an arithmetic pass-average of local
power-mel or log-mel maps; mel filtering and local-map averaging are not
interchangeable operations. Although originally
developed for audio, the mel-style filter bank is used here as a
logarithmic frequency-compression scheme rather than as a psychoacoustic
model. The mel-scale maps frequency $f$ (Hz) to
\begin{equation}
m = 2595 \log_{10}\!\left(1 + \frac{f}{700}\right),
\label{eq:mel}
\end{equation}
For the AE channel, the spectral representations cover the full Nyquist
band for empirical discrimination, but components above 1~MHz lie outside
the historically reported nominal AE response band and are interpreted
cautiously.
The full-resolution dB and 64-bin log-mel representations differ in both
frequency binning and amplitude transform. The dB-z caches differ from the
dB caches only by per-sample, per-plane z-standardisation. The canonical and
current RF spectral pipelines do not apply an additional fold-wise scaler to
spectrogram bins; therefore, the comparison is not a pure test of
mel-scale compression. The log-mel
variants (\texttt{ae\_mel}, \texttt{vib\_mel}) do not receive
per-sample normalisation; they therefore offer a partial sensitivity check
on whether removing absolute-energy information alters predictive
performance. A strict ablation omitting only per-sample/channel standardisation
is reported in Supplementary
Table~\ref{tab:supp-psnorm-ablation} and gives a small numerical mean-MAE
increase from 0.0210 to 0.0213~\(\upmu\)m without per-sample
standardisation. The maximum condition MAE increases from 0.1478 to
0.1863~\(\upmu\)m (Condition~7), so the ablation supports
mean-level robustness but not tail robustness.

Exact cache definitions, shapes, and SHA-256 provenance checksums are listed in Supplementary Table~\ref{tab:supp-cache-provenance}.
The available provenance for the two physical AE columns is separated from
the unavailable transfer-function information in Supplementary
Table~\ref{tab:supp-ae-plane-provenance}.

\subsubsection{Wavelet scattering and hand-crafted features}

Two-dimensional wavelet scattering transforms
(\texttt{ae\_wst}, \texttt{vib\_wst}) are computed with
\texttt{kymatio} \cite{mallat2012group}. We use second-order scattering
(\texttt{max\_order=2}) with $J=3$ octaves, Morlet wavelets, and
$L=8$ angular orientations (the constructor is
\texttt{Scattering2D(J=3, shape=(H,W), L=8, max\_order=2)}). Scattering coefficients are averaged over the
spatial (frequency--time) dimensions, producing a compact,
translation-invariant descriptor for each channel. The computation is
shown schematically in Figure~2. Note that for the vibration channel, the temporal dimension of the spectrogram is small ($T=13$ time frames). Using $J=3$ octaves corresponds to a maximum spatial support of $2^J = 8$ frames, which is close to the total signal length. Consequently, the wavelet filters at the largest scale ($j=2$) suffer from significant edge/boundary effects, and spatial averaging over this scale acts as a global smoother that discards almost all localized temporal variation, retaining only static spectral structures.

\begin{figure}[htbp]
\centering
\resizebox{\linewidth}{!}{%
\begin{tikzpicture}[
  node distance=1.4cm,
  box/.style={pub/representation, minimum width=2.6cm, minimum height=0.9cm},
  arrow/.style={pub/arrow}
]
\node[box] (spec) {2D spectrogram\\input};
\node[box, right=of spec] (psi1) {Morlet wavelet\\filter bank\\$L=8, J=3$};
\node[box, right=of psi1] (mod1) {Modulus\\$|\\cdot|$};
\node[box, right=of mod1] (psi2) {Second-order\\Morlet filters};
\node[box, right=of psi2] (mod2) {Modulus\\$|\\cdot|$};
\node[box, below=1.0cm of mod2] (avg) {Spatial\\averaging};
\node[box, left=of avg] (out) {Scattering\\coefficients};
\draw[arrow] (spec) -- (psi1);
\draw[arrow] (psi1) -- (mod1);
\draw[arrow] (mod1) -- (psi2);
\draw[arrow] (psi2) -- (mod2);
\draw[arrow] (mod2) -- (avg);
\draw[arrow] (avg) -- (out);
\end{tikzpicture}%
}
\caption{Wavelet scattering transform pipeline. A 2D spectrogram is
passed through two Morlet filter banks separated by modulus
nonlinearities, then spatially averaged to produce a compact,
translation-invariant descriptor per grinding pass.}
\label{fig:wst-pipeline}
\end{figure}

Time-domain statistics (\texttt{ae\_features}, \texttt{vib\_features})
include mean, standard deviation, skewness, kurtosis, RMS, peak, crest
factor, and frequency-domain band energy. Process parameters
(\texttt{pp}) are $v_s$, $v_w$, $a_p$. Physics-informed features
(\texttt{physics}) are heuristic indicators derived from process
parameters and sensor envelopes rather than dimensionally exact physical
quantities: AE wavelet energy (narrow 50--400\,kHz and broad 150--2000\,kHz
bands; the upper limit is capped at the AE Nyquist frequency of 2\,MHz,
and for the \texttt{cmor3-3} scales used the effective maximum frequency
is approximately 1.33\,MHz), AE burst rate (for each pass and band, $\sigma$ is computed from the
pass-level AE envelope after bandpass filtering; burst rate is the count
of envelope peaks above $4\sigma$, separated by at least 0.1\,ms. Because
$\sigma$ is computed per pass, the threshold is self-normalised and burst
rate captures relative burstiness rather than absolute AE amplitude), an equivalent chip-thickness proxy ($ec$), a normalized depth-of-cut
indicator ($bid=a_p/30$), a temperature-proxy term ($st$), envelope kurtosis for
the X/Y/Z vibration axes, and an overall vibration magnitude ($mag$). Each
indicator is aggregated over time by mean, standard deviation, maximum,
and minimum, producing a nominal 44-dimensional vector. A fold-wise audit
finds five globally constant aggregate columns (the two burst-rate minima
and the standard-deviation aggregates of (ec), (bid), and (st));
thus 39 columns vary across the campaign. These zero-variance columns are
retained for reproducibility, but fold-wise standardisation maps them to
zero and they must not be interpreted as independent physical
measurements. The audit is reported in Supplementary
Table~\ref{tab:supp-physics-audit}. The implemented process-derived and
temperature-indicator algorithms are specified directly in Supplementary
Table~\ref{tab:physics_features} and Equations~\eqref{eq:ec-implemented}--\eqref{eq:st-implemented}.
The archive does not preserve all filter settings for the burst features,
so those cached intermediate-derived values can be reused but are not
claimed to be fully reconstructible from raw signals. Because the
archived notes report the nominal AE response only up to 1~MHz, the 1--2\,MHz portion of
the broad band is retained only as an empirical descriptor and is not
interpreted as calibrated AE energy.

Tabular feature vectors and process parameters are standardised per fold
using training-set statistics. Spectrogram caches are not fold-wise scaled
in the canonical RF and current RF sensitivity pipelines. The dB-z cache is
instead normalised per sample and per channel before fitting; this removes
absolute energy information that may be physically meaningful. A strict
ablation that keeps the same RF architecture and omits only this per-sample
normalisation gives an overall LOGO
MAE of 0.0213~\(\upmu\)m, compared with 0.0210~\(\upmu\)m with
per-sample normalisation (Supplementary Table~\ref{tab:supp-psnorm-ablation}).
The result therefore does not
appear to depend on per-sample energy removal for mean LOGO MAE;
frequency-shape information appears sufficient for mean LOGO MAE in
this RF/dB-z configuration, although tail performance changes. Multimodal
tabular inputs are concatenated; multimodal deep-learning inputs are fed
through separate encoder branches that are fused before the regression
head.

\begin{figure}
\centering
\includegraphics[width=1\linewidth,height=\textheight,keepaspectratio]{methods_signal_representations.png}
\caption{Example signal representations for condition 10, sample 1
($R_a = 0.098~\upmu$m). (a)~AE dB spectrogram, (b)~AE log-mel
spectrogram, (c)~vibration (Z-axis) dB spectrogram, and
(d)~vibration log-mel spectrogram. The examples show dominant lower-frequency
energy in both modalities. Vibration below 5~kHz is compatible with structural
dynamics, whereas AE is treated as an empirical contact-related descriptor
because its transfer function is unavailable. Higher-frequency AE energy above 0.5~MHz and vibration energy near
17~kHz are located in frequency ranges commonly associated with
grain-contact and chatter-related dynamics, respectively.}
\label{fig:representations}
\end{figure}

\subsection{Predictive models}\label{predictive-models}

The model families share the outer LOGO test folds and metrics, with
model-specific training, validation, and preprocessing procedures
(Figure~\ref{fig:taxonomy}).

\textbf{Auditable classical baselines.} Ridge coefficients can be inspected
directly; the tree ensembles are analysed with post-hoc attribution. This is relevant to monitoring
workflows where model updates may be required after changes in wheel
grade, workpiece material, or dressing condition, although this workflow
is not tested here.

\begin{itemize}
\tightlist
\item
  \textbf{Ridge regression:} L2-regularised linear regression with
  $\alpha=1.0$; coefficients provide a signed importance score for every
  frequency bin.
\item
  \textbf{Random forest:} 200 trees, maximum depth 8, minimum 1 sample
  per leaf, $n_{\text{jobs}}=1$.
\item
  \textbf{LightGBM and XGBoost:} gradient-boosted trees with 200
  estimators, maximum depth 5, learning rate $0.05$, subsample
  $0.8$, and column sampling $0.8$.
\item
  \textbf{Shallow MLP:} two hidden layers of 128 and 64 units, L2
  penalty $\alpha=10^{-4}$, Adam optimiser with learning rate
  $10^{-3}$, early stopping on 10\% validation data, maximum 500
  iterations.
\end{itemize}

\textbf{Deep-learning baselines.} These models represent recent
literature on grinding monitoring:

\begin{itemize}
\tightlist
\item
  ResNet-style encoders for AE and vibration spectrograms
  (ResNetAECNN, ResNetVibCNN).
\item
  ResNet-based multimodal fusion (ResNetFusion).
\item
  Bilinear attention-based fusion network (BilinearFusionNetwork).
\item
  Multi-resolution spectrogram CNN.
\item
  Trajectory-based sequence CNN (TrajectoryCNN).
\item
  Graph-neural-network fusion (GraphNeuralNetworkFusion), TabTransformer (TabTransformerRegressor), and attention-based
  regressors.
\end{itemize}

Specifically, the advanced deep-learning baselines are structured as follows:
\begin{itemize}
\item \textbf{BilinearFusionNetwork:} uses dual ResNet encoders to extract feature maps from AE and vibration spectrograms, then applies low-rank bilinear pooling to fuse these representations with process parameters and physics-informed indicators before MLP regression.
\item \textbf{TrajectoryCNN:} is a sequence-based model that utilizes 1D temporal convolutions to process temporal sub-band energy trajectories from AE and vibration, capturing localized time-dependent dynamics instead of relying on time-averaged summaries.
\item \textbf{GraphNeuralNetworkFusion:} frames multimodal fusion as a graph learning task where individual modalities (AE, vibration, physics, parameters) represent nodes in a fully-connected graph. Node features are processed via Graph Attention Network (GAT) layers for cross-modal message passing, followed by global average graph pooling.
\item \textbf{TabTransformerRegressor:} embeds continuous features (physics, process parameters) by scaling a learned feature-specific embedding vector by the continuous scalar value and adding a bias. These embedded feature sequences are then processed by a multi-layer Transformer encoder.
\end{itemize}

The best-performing spectrogram CNN, ResNetVibCNN, processes vibration
dB spectrograms of shape $(3, 257, 13)$ through a stem convolution,
two residual blocks with squeeze-and-excitation attention and dropout
$p=0.3$, global average pooling, and a fully connected regressor with a
64-dimensional hidden layer. Its architecture is shown in
Figure~4. All deep-learning models are trained
with AdamW (learning rate $10^{-3}$, weight decay $10^{-4}$), MSE loss,
a ReduceLROnPlateau scheduler (factor $0.5$, patience 5 epochs), early
stopping (patience 20 epochs), and a maximum of 200 epochs. Batch sizes
are 16 for the ResNet encoders and 32 otherwise. The historical planned
comparison uses the canonical seed-42 results for all model--configuration
pairs; reproducible current RF reruns are reported before those artifacts in
the Results. To quantify training stochasticity, a secondary equal-budget
sensitivity analysis repeats the main stochastic architectures over five
seeds; details are provided in Section~\ref{experimental-workflow}.

\begin{figure}[htbp]
\centering
\resizebox{0.98\linewidth}{!}{\begin{tikzpicture}[
    node distance=0.72cm and 0.82cm,
    stage/.style={pub/model, text width=4.25cm, minimum height=1.12cm},
    arrow/.style={pub/arrow},
]
\node[pub/input, text width=4.25cm, minimum height=1.12cm] (input)
  {Vib-dB input\\$3\times257\times13$};
\node[stage, right=of input] (stem)
  {Stem\\$3\times3$ Conv + BN + ReLU\\$16\times257\times13$};
\node[stage, right=of stem] (res1)
  {SE-residual block 1\\$3\times3$ Conv $\times2$, stride 2\\$32\times129\times7$};

\node[stage, below=of res1] (res2)
  {SE-residual block 2\\$3\times3$ Conv $\times2$, stride 2\\$64\times65\times4$};
\node[pub/explanation, text width=4.25cm, minimum height=1.12cm, left=of res2] (gap)
  {Global average pooling\\$64$ features};
\node[pub/target, text width=4.25cm, minimum height=1.12cm, left=of gap] (head)
  {Dense head and output\\$64$ ReLU + dropout $p=0.3$\\Measured $R_a$ (\textmu m)};

\draw[arrow] (input) -- (stem);
\draw[arrow] (stem) -- (res1);
\draw[arrow] (res1) -- (res2);
\draw[arrow] (res2) -- (gap);
\draw[arrow] (gap) -- (head);
\node[pub/note, text width=13.6cm, align=center, below=0.55cm of gap]
  {Each residual block applies batch normalization, squeeze-and-excitation attention, and dropout $p=0.3$.};
\end{tikzpicture}%
}
\caption{ResNetVibCNN architecture for vibration dB spectrogram
regression. The three-axis input is processed by a stem convolution and
two stride-two SE-residual blocks. Each residual block contains batch
normalisation, squeeze-and-excitation attention, and dropout; global average
pooling and a 64-unit dense head produce the measured-$R_a$ estimate.}
\label{fig:resnetvibcnn-arch}
\end{figure}

Table~\ref{tbl:dl-baselines} summarises the deep-learning baselines and
their input configurations.

\begin{table}[htbp]
\centering
\caption{Overview of the deep-learning baselines. The column ``Input''
uses the abbreviations defined in Table~\ref{tbl:config-abbreviations}.}
\label{tbl:dl-baselines}
\begin{tabular}{@{}llll@{}}
\toprule
Model & Input & Key mechanism & Fusion \\
\midrule
ResNetAECNN & AE-dB & ResNet encoder + SE & Single modality \\
ResNetVibCNN & Vib-dB & ResNet encoder + SE & Single modality \\
ResNetFusion & AE-dB + Vib-dB & Dual ResNet encoders & Early fusion \\
BilinearFusion & AE-dB + Vib-dB + phys + PP & Bilinear attention & Late fusion \\
Multi-resolution CNN & Vib-dB & Multi-scale convolutions & Single modality \\
Trajectory CNN & \texttt{Vib-traj} & 1D sequence CNN on temporal sub-bands & Single modality \\
GNN fusion & AE-dB + Vib-dB & Graph neural network & Structured fusion \\
TabTransformer & \texttt{AE/Vib feats + phys + PP} & Transformer on tabular & Tabular multimodal \\
Attention regressor & AE-dB + Vib-dB & Self-attention & Late fusion \\
\bottomrule
\end{tabular}
\end{table}

\begin{figure}
\centering
\resizebox{0.98\linewidth}{!}{\begin{tikzpicture}[
    arrow/.style={pub/arrow},
    family/.style={pub/explanation, text width=12.4cm, minimum height=0.82cm},
    modelcard/.style={pub/model, text width=2.65cm, minimum height=0.88cm},
]
\node[pub/input, text width=11.6cm, minimum height=0.88cm] (input) at (0,6.4)
  {Input representations: AE, vibration, process parameters, and physics-inspired features};
\node[pub/heading] (choice) at (0,5.45)
  {Alternative model families evaluated under the same LOGO protocol};
\draw[arrow] (input) -- (choice);

\node[family] (auditable) at (0,4.45)
  {Auditable baselines: direct coefficients, split structure, and post-hoc TreeSHAP};
\draw[arrow] (choice) -- (auditable);
\node[modelcard] (ridge) at (-5.0,3.25) {Ridge regression};
\node[modelcard] (rf) at (-1.7,3.25) {Random forest};
\node[modelcard] (lgbm) at (1.7,3.25) {LightGBM /\\XGBoost};
\node[modelcard] (mlp) at (5.0,3.25) {Shallow MLP};
\node[pub/note, text width=12.0cm, align=center] at (0,2.35)
  {Ridge coefficients, tree split structure, and TreeSHAP support model inspection; ensemble models remain complex fitted predictors.};

\node[family] (deep) at (0,1.35)
  {Deep-learning baselines: learned spectral and multimodal predictors with Grad-CAM and MC-dropout diagnostics};
\node[modelcard] at (-5.0,0.15) {ResNet CNNs\\(AE / vibration)};
\node[modelcard] at (-1.7,0.15) {ResNet fusion};
\node[modelcard] at (1.7,0.15) {Bilinear fusion};
\node[modelcard] at (5.0,0.15) {Multiscale CNN};
\node[modelcard] at (-3.35,-1.05) {Trajectory CNN};
\node[modelcard] at (0,-1.05) {GNN /\\TabTransformer};
\node[modelcard] at (3.35,-1.05) {Attention regressor};

\node[pub/note, text width=12.0cm, align=center] at (0,-1.95)
  {Every listed configuration regresses the same averaged measured target.};
\node[pub/target, text width=4.0cm, minimum height=0.88cm] at (0,-2.8)
  {Predict measured $R_a$ (\textmu m)};
\end{tikzpicture}%
}
\caption{Model taxonomy evaluated in this study. The left branch contains
a directly inspectable regularised linear model and auditable tree ensembles
supported by post-hoc SHAP values. The right branch uses deep-learning
baselines whose decisions are analysed post-hoc with Grad-CAM and
MC-dropout. Both branches predict $R_a$ and are compared under the same
LOGO protocol.}
\label{fig:taxonomy}
\end{figure}

\begin{table}[htbp]
\centering
\caption{Key hyperparameters of the main auditable classical and baseline models.
Fixed hyperparameters were applied uniformly across all LOGO folds; deep-learning
entries (Shallow MLP) used additional early-stopping on the validation fold.}
\label{tbl:hyperparams}
\setlength{\tabcolsep}{8pt}
\renewcommand{\arraystretch}{1.25}
\begin{tabular}{@{}llr@{}}
\toprule
\textbf{Model} & \textbf{Hyperparameter} & \textbf{Value} \\
\midrule
Ridge                & Regularisation strength $\alpha$   & 1.0 \\
\addlinespace[4pt]
Random forest        & Number of trees                    & 200 \\
                     & Max depth                          & 8 \\
                     & Min samples per leaf               & 1 \\
\addlinespace[4pt]
LightGBM             & Number of estimators               & 200 \\
                     & Max tree depth                     & 5 \\
                     & Learning rate                      & 0.05 \\
                     & Subsample ratio                    & 0.8 \\
\addlinespace[4pt]
XGBoost              & Number of estimators               & 200 \\
                     & Max tree depth                     & 5 \\
                     & Learning rate                      & 0.05 \\
                     & Subsample ratio                    & 0.8 \\
\addlinespace[4pt]
Shallow MLP          & Hidden layers                      & (128, 64) \\
                     & L2 regularisation $\alpha$         & $10^{-4}$ \\
                     & Learning rate                      & $10^{-3}$ \\
\bottomrule
\end{tabular}
\end{table}

\subsection{Validation protocol}\label{validation-protocol}

We use strict leave-one-condition-out (LOGO) cross-validation with 16
outer folds. In each fold, one condition is held out as the test set, a
second condition is reserved for early-stopping validation, and the
remaining 14 conditions are used for training
(Figure~\ref{fig:logo}). The validation condition is selected from the
remaining conditions by successive draws from
\texttt{numpy.random.RandomState(42)} in ascending test-condition order.
The resulting train/validation/test split is
therefore identical for every model evaluated on a given fold. This prevents
information leakage across samples that share the same grinding
condition and measures whether a model generalises to previously unseen
process states. Typical fold sizes are approximately 280 training,
20 validation, and 20 test samples. The deterministic fold assignment is
shown in Figure~\ref{fig:logo-matrix}. Because only one validation
condition is used per fold, the validation loss is potentially volatile and
sensitive to that condition's characteristics. A fixed-hyperparameter
rotation diagnostic is reported for ResNetVibCNN and ResNetAECNN in
Supplementary Table~\ref{tab:supp-deep-validation-rotation}: it compares
the canonical seed-42 validation draw with a cyclic-next-condition
validation assignment. It tests early-stopping sensitivity only; it does
not rerun the original hyperparameter search or provide inner grouped model
selection. Tabular models such as RF use fixed hyperparameters and do not
rely on this early-stopping signal.
The exact 16 assignments are listed in Supplementary
Table~\ref{tab:supp-logo-assignments}; Figure~\ref{fig:logo-matrix} is
generated from the same rule and archived assignment CSV.

A model-selection caveat applies. The deep-learning entries were tuned
by random search over model and training hyperparameters (with a hyperparameter search space size of approximately 20 configurations per architecture) using the same
outer LOGO folds (the validation fold was used for early stopping and the
test fold was used for final model selection). The tree-based models used
fixed hyperparameters (Table~\ref{tbl:hyperparams}). This is not a fully
nested leave-one-condition-out comparison; using the held-out test fold
for final model selection means the deep-learning results are
post-selection estimates rather than unbiased generalisation estimates.
The reported differences should therefore be read as rankings under a
fixed evaluation budget rather than as population estimates of
generalisation error. This post-selection caveat applies to the selected
best RF as well as to the deep-learning entries. The primary purpose of
the benchmark is comparative evidence within this dataset, not an
unbiased population estimate.

The primary planned statistical testing is performed on the canonical seed-42 results using paired Wilcoxon signed-rank tests. However, because training sets overlap substantially across folds (each pair of training sets shares 13 out of 14 conditions), the fold MAEs are not strictly independent. This violation of the independence assumption means the resulting p-values must be interpreted as exploratory indicators of relative performance rather than strict, independent probability assessments. A secondary sensitivity analysis using five random seeds (42--46) is conducted to assess training stochasticity, as described in Section~\ref{experimental-workflow}.

\begin{figure}[htbp]
\centering
\includegraphics[width=89mm]{methods_logo_matrix.png}
\caption{Leave-one-condition-out fold assignment. Rows correspond to the
16 LOGO folds; columns correspond to the 16 grinding conditions. In each
fold the test condition is red, the validation condition is orange, and
the 14 training conditions are blue.}
\label{fig:logo-matrix}
\end{figure}

\begin{figure}
\centering
\begin{tikzpicture}[
    train/.style={pub/train},
    val/.style={pub/validation},
    test/.style={pub/test},
    arrow/.style={pub/arrow},
    ann/.style={pub/heading},
]

% One representative LOGO fold. The caption supplies the full figure title.
\foreach \c in {1,...,16} {
    \pgfmathsetmacro{\col}{mod(\c-1,8)}
    \pgfmathsetmacro{\row}{floor((\c-1)/8)}
    \pgfmathsetmacro{\x}{\col*1.10}
    \pgfmathsetmacro{\y}{0.45 - \row*1.12}
    \ifnum\c=1
        \node[test] (C\c) at (\x,\y) {C\c\\TEST};
    \else\ifnum\c=8
        \node[val] (C\c) at (\x,\y) {C\c\\VAL};
    \else
        \node[train] (C\c) at (\x,\y) {C\c};
    \fi\fi
}

\node[ann] (annTest) at (0.0,-2.35) {Held-out test condition};
\draw[arrow] (annTest.north) -- (C1.south);
\node[ann] (annVal) at (4.7,-2.35) {Validation condition};
\draw[arrow] (annVal.north) -- (C8.south);

\node[pub/base, fill=white, text width=7.5cm] at (3.85,-3.65) {
  The remaining 14 conditions train the model. Repeat the split 16 times so each condition is tested once.
};

\node[train, minimum width=0.8cm, minimum height=0.45cm] (legT) at (0.2,-4.8) {};
\node[right=0.08cm of legT, pub/legend] {Train (14)};
\node[val, minimum width=0.8cm, minimum height=0.45cm] (legV) at (3.0,-4.8) {};
\node[right=0.08cm of legV, pub/legend] {Validation (1)};
\node[test, minimum width=0.8cm, minimum height=0.45cm] (legTe) at (6.3,-4.8) {};
\node[right=0.08cm of legTe, pub/legend] {Test (1)};
\end{tikzpicture}%

\caption{Leave-one-condition-out cross-validation, illustrated with
canonical fold 1 (Condition 1 test; Condition 8 validation). Each of the 16
conditions is used once as the test fold (red); a different condition
serves as the validation fold (orange); the remaining 14 conditions are
used for training (blue). Per-fold performance is summarised by MAE, and
the 16 per-fold MAEs are aggregated as mean~$\pm$~std. Pairwise model
comparisons use the Wilcoxon signed-rank test with Holm-Bonferroni
correction.}
\label{fig:logo}
\end{figure}

The primary metric is mean absolute error (MAE) in micrometres:
\begin{equation}
\mathrm{MAE} = \frac{1}{N}\sum_{i=1}^{N} |y_i - \hat{y}_i|.
\label{eq:mae}
\end{equation}
We report mean~$\pm$~standard deviation of MAE across the 16 folds.
Bootstrap 95\% confidence intervals are computed with $B=1000$
replicates on the per-fold MAEs using the percentile method. Pairwise
comparisons use the two-sided Wilcoxon signed-rank test paired by LOGO
fold, with Holm-Bonferroni correction at family-wise error rate
$\alpha=0.05$. This statistical workflow is shown in
Figure~\ref{fig:statistical-workflow}.
The implementation uses \texttt{scipy.stats.wilcoxon} with
\texttt{zero\_method='zsplit'}; tied absolute differences receive average
ranks. Hodges--Lehmann intervals use 5,000 percentile bootstrap resamples
of the 16 paired fold differences with random seed 42. The point estimate
is the median of all Walsh averages $\{(d_i+d_j)/2:i\leq j\}$ of the
paired differences. These intervals are unadjusted descriptive intervals;
multiplicity correction applies only to the Wilcoxon $p$-values. Rank-biserial
correlation is the signed-rank sum divided by the total rank sum after
removing zero differences. No exact zero paired differences occur in the
reported 14 comparisons, so this removal does not create a discrepancy with
the Wilcoxon zero-handling rule.

\begin{figure}[htbp]
\centering
\begin{tikzpicture}[
  node distance=1.6cm,
  box/.style={pub/representation, minimum width=3.2cm, minimum height=0.9cm},
  arrow/.style={pub/arrow}
]
\node[box] (mae) {16 per-fold MAEs\\for each model};
\node[box, right=of mae] (boot) {Bootstrap\\$B=1000$ replicates\\(percentile CI)};
\node[box, below=of mae] (pair) {Pair by fold\\Wilcoxon signed-rank};
\node[box, below=of pair] (holm) {Holm--Bonferroni\\correction};
\node[box, right=of holm] (decide) {Reject/retain\\at $\alpha=0.05$};
\draw[arrow] (mae) -- (boot);
\draw[arrow] (mae) -- (pair);
\draw[arrow] (pair) -- (holm);
\draw[arrow] (holm) -- (decide);
\end{tikzpicture}
\caption{Statistical comparison workflow. Per-fold MAEs are used both to
build bootstrap confidence intervals and to perform paired Wilcoxon tests;
$p$-values are adjusted with Holm--Bonferroni correction.}
\label{fig:statistical-workflow}
\end{figure}

\subsection{Explainability and uncertainty}\label{explainability-and-uncertainty}

To assess whether the predictions can be audited, selected top models
are analysed using coefficient-, SHAP-, Grad-CAM-, and uncertainty-based
diagnostics.

\textbf{Model-attribution analyses.}
Ridge regression yields a signed coefficient for every frequency bin,
providing a global linear attribution; Figure~\ref{fig:xai-methods}a
shows the coefficient magnitudes for \texttt{vib\_spec}, which
concentrate in the low-frequency chatter and structural bands. TreeSHAP is applied to both (i)~the LightGBM model trained on
\texttt{ae\_spec+vib\_spec} as a representative tree-based model, and
(ii)~the recommended random forest trained on fused AE and vibration
dB-z descriptors (\texttt{AE-dB-z + Vib-dB-z}). For the
LightGBM model the training set is used as the background distribution
with exact TreeSHAP computations; for the random forest we use the
path-dependent TreeSHAP approximation. The LightGBM vibration profile
contains a sparse high-importance bin near 22~kHz. This lies in a
structural-frequency region, but neither a specific grinding mechanism nor
sensor/fixture resonance can be separated from the present data
(Figure~\ref{fig:xai-methods}b). The full ranked
SHAP importance for the LightGBM model is shown in Figure~\ref{fig:xai},
and the pooled out-of-fold SHAP profile for the current 14-condition RF/dB-z
pipeline is shown in Supplementary Figure~\ref{fig:supp-shap-rf}. It shows
AE importance above approximately
300~kHz and dominant vibration importance between approximately 16 and
24~kHz, with the largest pooled bin near 16.2~kHz. The profiles share a broad separation between
higher-frequency AE information and structural-frequency vibration
information, but their dominant bins differ by model and representation.
For each TreeSHAP analysis, feature importance is first defined as the mean
absolute SHAP value over the sampled explanatory passes. It is then averaged
over STFT time frames, retained separately by AE plane or vibration axis,
and summed within each reported physical-frequency band. Each modality and
explainer is normalised independently to 100\% only after this pooling.
The flattened arrays use NumPy C order. Within a modality, flattened feature
index $k=((cF)+f)T+t$ maps to channel/plane $c$, frequency bin $f$, and time
frame $t$; the vibration block begins after the $2\times300\times47$ AE
features. Physical frequency is $f f_s/N_{\rm FFT}$. This mapping is used by
both the plotted profiles and the archived per-frequency CSV files.
The attribution protocols are distinct. First, LightGBM uses a global
explanatory refit with 30 background and 50 explained passes as an
illustrative exact-tree explanation. Second, the RF analysis refits the
reproducible current 14-condition RF/dB-z pipeline in each LOGO fold and
explains only that fold's held-out condition using path-dependent TreeSHAP;
fold-specific profiles are pooled only after out-of-fold explanation. It
therefore explains the current RF predictions, not the historical canonical
0.0200~\(\upmu\)m artifact. Third, ResNetVibCNN uses Grad-CAM as described
below. No global RF explanation is used for the reported RF frequency profile.

For the deep-learning baseline ResNetVibCNN, Grad-CAM is computed on the
last convolutional layer to visualise the spectrogram regions that drive
predictions. The aggregated Grad-CAM frequency profile highlights chatter
and structural frequencies below 15~kHz, with additional high-frequency
vibration sensitivity above 15~kHz that may reflect wheel--spindle
structural response or high-frequency vibration modes. Because the input
is vibration, this high-frequency sensitivity should not be interpreted as
acoustic-emission energy (Figure~\ref{fig:xai-methods}c).
For Grad-CAM, the absolute activation map is averaged over all analysed
passes after bilinear interpolation from the final convolutional map to the
original \(257\times13\) vibration-STFT grid, then averaged over time frames
to produce the frequency profile. Thus CSV frequency bin \(k\) maps to
\(k\times 51{,}200/512\)~Hz. Interpolation supplies a common display grid;
the saliency still represents the receptive fields of the pooled
convolutional features rather than an isolated original-STFT bin.

\begin{figure}[htbp]
\centering
\includegraphics[width=\textwidth]{methods_xai_comparison.png}
\caption{Explainability methods applied to real grinding data. (a)~Ridge
regression coefficient magnitudes for Vib-dB. (b)~TreeSHAP
mean $|\text{SHAP}|$ for LightGBM on Vib-dB.
(c)~Grad-CAM frequency importance for ResNetVibCNN on Vib-dB.
Shaded bands mark low-frequency structural modes (green) and
higher-frequency vibration/chatter regions (blue).}
\label{fig:xai-methods}
\end{figure}

\textbf{Physics-consistency checks.}
Important bins are mapped to physical frequencies. For AE spectrograms
the bin spacing is approximately $6.69$~kHz/bin; for vibration
spectrograms it is approximately $100$~Hz/bin. A bin is considered
consistent with a known grinding mechanism if its centre frequency lies
within $\pm$10\% of an expected peak frequency for that mechanism; the
tolerance is intended only for narrow peak matching, whereas broad bands
such as 2--15~kHz are interpreted qualitatively. The band edges are
approximate and the tolerance is applied to the centre frequency rather
than to the band edges. In this study the expected bands are approximated
as: spindle rigid-body motion and low-order spindle drift below 500~Hz;
chatter and structural vibration roughly 2--15~kHz; and grain-contact
acoustic emission above 100~kHz, with the historically reported
100~kHz--1~MHz nominal band used only as a cautious reporting convention.
Figure~\ref{fig:bandmass} shows how the total attribution
mass is distributed across physical frequency bands for the three
explainers; the mass is concentrated in frequency regions that are
physically plausible for grinding dynamics.

\textbf{Uncertainty quantification.}
MC-dropout with 50 stochastic forward passes is used to quantify
predictive uncertainty for ResNetVibCNN on Vib-dB. All 16
uncertainty-analysis LOGO checkpoints are loaded without retraining,
substitution, or non-strict state-dictionary matching. Deterministic and
stochastic predictions are recomputed from each same checkpoint in one run.
These checkpoints give a deterministic mean fold MAE of 0.0443~\(\upmu\)m
and are not the distinct canonical 0.0268~\(\upmu\)m result artifact.
Dropout layers remain active at test time while
batch-normalisation layers are kept in evaluation mode so that small test
batches do not shift the population statistics. All dropout layers in the
network are activated for the stochastic passes, so the stochastic-forward
mean is treated as an uncertainty diagnostic and not as a recalibrated
deterministic predictor. The mean and variance of the 50 outputs form the
point prediction and the epistemic standard deviation, respectively. The
reported 95\% interval is computed as
\(\hat y \pm 1.96\,\sigma_{\text{MC}}\) and should be interpreted as an
illustrative epistemic interval rather than a calibrated predictive
interval including aleatoric uncertainty. The MC-dropout analysis is
evaluated under the same 16-fold LOGO protocol used for accuracy, so
every condition---including Condition~7---appears as a held-out test set
once. Figure~\ref{fig:uncertainty-methods} shows the resulting prediction
intervals sorted by predicted mean, a reliability diagram relating
predicted standard deviation to observed absolute error, and the
empirical coverage within each uncertainty bin. Across all 319 held-out
samples the aggregate coverage is 50.2\%, the mean interval width is
0.0479~\(\upmu\)m, the mean absolute error of the stochastic-mean prediction
is 0.0452~\(\upmu\)m, and the Pearson correlation between predicted standard
deviation and observed absolute error is 0.488. Although predicted
standard deviation is therefore correlated with absolute error, the
intervals remain far below the nominal 95\% level; MC-dropout does not
provide a calibrated predictive interval for this CNN. Condition~7 has
0\% coverage. A smaller 59-sample subset (Conditions 1, 2, and 6)
previously gave 100\% coverage but with a wider mean width
(0.073~\(\upmu\)m versus 0.0479~\(\upmu\)m under full LOGO), showing that the
apparent calibration of MC-dropout is highly sensitive to which conditions
are included.
Figure~\ref{fig:uncertainty-calibration} shows the per-sample
relationship between predicted standard deviation and observed absolute
error for the full LOGO evaluation.

For a separate 15-condition RF uncertainty variant, intervals are computed as
\(\hat y \pm c \, \sigma_{\text{tree}}\), where
\(\sigma_{\text{tree}}\) is the standard deviation of the individual
tree predictions. Within outer fold \(g\), let \(\bar y_i^{\rm OOB}\) and
\(s_i^{\rm OOB}\) be the mean and standard deviation over trees for which
training sample \(i\) is out of bag. The global factor is
\[
c_g=Q_{0.95}\left\{\frac{|y_i-\bar y_i^{\rm OOB}|}
{s_i^{\rm OOB}+10^{-8}}:i\in\mathcal T_g\right\},
\]
where \(Q_{0.95}\) is NumPy's default linear empirical quantile. The
condition-aware variant uses the same score restricted to samples from the
three Euclidean-nearest training conditions after min--max normalisation of
wheel speed, workpiece speed, and depth. The split variant removes the
single nearest condition from model fitting and applies the same scaled-score
quantile on that calibration condition. The naive variant uses
\(Q_{0.95}(|y_i-\bar y_i^{\rm OOB}|)\) directly, without tree-variance
scaling. Zero tree standard deviations are stabilised by \(10^{-8}\), and
all fitting, neighbour selection, and calibration are repeated independently
inside each outer fold. These are diagnostic post-hoc intervals, not formal
finite-sample conformal guarantees or intervals for the current 14-condition
benchmark RF.

\begin{figure}[htbp]
\centering
\includegraphics[width=\textwidth]{methods_uncertainty_overview.png}
\caption{Full 16-fold LOGO MC-dropout uncertainty for ResNetVibCNN on
Vib-dB. (a)~Nominal 95\% MC-dropout epistemic intervals and observed $R_a$
sorted by predicted mean. (b)~Reliability diagram: mean predicted
standard deviation versus mean observed absolute error per bin.
(c)~Empirical coverage per uncertainty bin against the nominal 95\%
level. Aggregate coverage is 50.2\% and mean interval width is
0.0479~\(\upmu\)m, indicating that MC-dropout is under-covering on this
CNN; the intervals are correlated with error but are not calibrated.}
\label{fig:uncertainty-methods}
\end{figure}

\subsection{Experimental workflow}\label{experimental-workflow}

The study follows a four-stage workflow: (1)~preprocessing of raw AE and
vibration signals and caching of all candidate representations; (2)~LOGO
training and evaluation of every model--configuration pair under the
same protocol; (3)~ranking, bootstrap confidence intervals, and
statistical testing; and (4)~post-hoc explainability and uncertainty
analysis of the top-performing systems. A high-level overview of the
pipeline is shown in Figure~\ref{fig:workflow}.

Because the 16 LOGO folds are deterministic and exhaustive over the
available grinding conditions, repeating the outer cross-validation would
not create new validation evidence. The historical planned comparison uses
the canonical seed-42 results and the pre-defined 14-comparison
Holm-corrected testing family.  To quantify training stochasticity, a
secondary sensitivity analysis trained the main stochastic architectures
(ResNetVibCNN, TrajectoryCNN, ResNetAECNN, ChannelAttentionCNN, ResNetFusion, BilinearFusionNetwork, and
the recommended random forest) under the same 16 outer LOGO folds with
five different random seeds. In every outer fold, the repeated RF excludes
the same canonical validation condition as the neural models and therefore
uses the same 14-condition training-data budget. The earlier 15-condition RF
repeat is retained only as a deployment-style sensitivity and is excluded
from the paired model-family tests. The reported seed-averaged MAE is the mean
over these five seeds; seed-level overall MAEs are given in
Supplementary Table~\ref{tab:supp-seed-level}.  For the sensitivity
analysis, per-condition MAEs were first averaged over seeds, preserving
the 16 condition-level paired observations; the 16 condition-level values
were not inflated by treating folds and seeds as independent samples.
This seed-repeat analysis provides a robustness check for training
stochasticity and is not treated as the primary inferential comparison.

The historical shallow-MLP baseline used scikit-learn's internal random
10\% validation fraction and is therefore an implementation-specific member
of the archived 14-comparison family. A corrected sensitivity uses the same
designated whole validation condition as the other models, standardises inputs
and targets from the 14-condition fitting subset only, disables the internal
random validation split, and selects among three training-only
architecture/regularisation settings. This corrected run is reported as a
separate sensitivity rather than substituted into the historical family.

\begin{figure}[htbp]
\centering
\resizebox{\linewidth}{!}{%
\begin{tikzpicture}[
    raw/.style={pub/input, minimum width=2.6cm, minimum height=0.85cm, text width=2.5cm},
    rep/.style={pub/representation, minimum width=2.6cm, minimum height=0.85cm, text width=2.5cm},
    pp/.style={pub/process, minimum width=2.6cm, minimum height=0.85cm, text width=2.5cm},
    phys/.style={pub/physics, minimum width=2.6cm, minimum height=0.85cm, text width=2.5cm},
    config/.style={pub/explanation, minimum width=4.0cm, minimum height=1.05cm, text width=3.8cm},
    model/.style={pub/model, minimum width=3.6cm, minimum height=1.05cm, text width=3.4cm},
    target/.style={pub/target, minimum width=2.8cm, minimum height=1.6cm, text width=2.6cm},
    arrow/.style={pub/arrow},
    flowarrow/.style={pub/flow},
    label/.style={pub/heading},
]

% ===== Left column: input modalities =====
\node[raw]  (ae_raw)  at (0, 5.0) {AE Raw\\(4.0 MHz)};
\node[rep]  (ae_spec) at (0, 3.5) {AE Spectrogram\\(2 channels$\times$300$\times$47)};
\node[rep]  (ae_feat) at (0, 2.0) {AE Features\\(RMS, Kurt, Skew, Crest)};
\draw[arrow] (ae_raw)  -- (ae_spec);
\draw[arrow] (ae_spec) -- (ae_feat);

\node[raw]  (vib_raw)  at (0, 0.0)  {Vib Raw\\(51.2 kHz)};
\node[rep]  (vib_spec) at (0,-1.5)  {Vib Spectrogram\\(3$\times$257$\times$13)};
\node[rep]  (vib_feat) at (0,-3.0)  {Vib Features\\(RMS, Kurt, Skew, Crest)};
\draw[arrow] (vib_raw)  -- (vib_spec);
\draw[arrow] (vib_spec) -- (vib_feat);

\node[pp]   (pp)   at (0,-5.0) {Process Params\\($v_s,\,v_w,\,a_p$)};
\node[phys] (phys) at (0,-6.5) {Physics Features\\(11 indic.$\times$4 agg. = 44-D)};

% Column headers
\node[label] at (0,   7.2) {Input Modalities};
\node[label] at (5.5, 7.2) {Input Configuration\\Categories};
\node[label] at (10.3,7.2) {Representative Model Families};

% ===== Flow arrows between columns =====
\draw[flowarrow] (1.5, 0.25) -- (3.3, 0.25) node[midway, above, pub/heading] {combine};
\draw[flowarrow] (7.7, 0.25) -- (8.7, 0.25) node[midway, above, pub/heading] {evaluate};
\draw[flowarrow] (12.1,0.25) -- (13.1,0.25) node[midway, above, pub/heading] {predict};

% ===== Middle column: configuration categories =====
\node[config] (c1) at (5.5, 5.6) {1.\ Process params only};
\node[config] (c2) at (5.5, 4.2) {2.\ Single modality\\(AE or Vib spectrum/features)};
\node[config] (c3) at (5.5, 2.8) {3.\ AE+Vib fusion\\(no physics)};
\node[config] (c4) at (5.5, 1.4) {4.\ AE+Vib+physics\\(+ process params)};
\node[config] (c5) at (5.5, 0.0) {5.\ Single modality\\+ physics features};
\node[config] (c6) at (5.5,-1.4) {6.\ Process params\\+ physics features};
\node[config] (c7) at (5.5,-2.8) {7.\ All modalities combined};

% ===== Right column: model zoo =====
\node[model] (m1) at (10.3, 5.6) {Classical ML\\Ridge / RF / XGB / LightGBM\\Shallow MLP};
\node[model] (m2) at (10.3, 4.2) {Single-Modality CNNs\\ResNetAE / ResNetVib\\ChannelAttention};
\node[model] (m3) at (10.3, 2.8) {Fusion CNNs\\Bilinear / ResNetFusion};
\node[model] (m4) at (10.3, 1.4) {Graph Fusion\\GNN};
\node[model] (m5) at (10.3, 0.0) {Tabular DL\\TabTransformer};
\node[model] (m6) at (10.3,-1.4) {Attention Fusion\\Cross-modal attention};
\node[model] (m7) at (10.3,-2.8) {Sequence Models\\TrajectoryCNN};

% ===== Target =====
\node[target] (target) at (14.5, 0.25) {Surface Roughness $R_a$\\(\textmu m)\\0.066--0.381};

% Representative dashed connectors
\draw[pub/dashed] (ae_spec.east)  -- ++(0.4,0) |- (c2.west);
\draw[pub/dashed] (ae_feat.east)  -- ++(0.4,0) |- (c7.west);
\draw[pub/dashed] (vib_spec.east) -- ++(0.4,0) |- (c3.west);
\draw[pub/dashed] (pp.east)       -- ++(0.4,0) |- (c1.west);
\draw[pub/dashed] (phys.east)     -- ++(0.4,0) |- (c6.west);

% ===== Compact two-row legend =====
\node[raw, minimum width=0.55cm, minimum height=0.36cm, text width=0.55cm] (leg1) at (1.0,-8.25) {};
\node[right=0.08cm of leg1, pub/legend] {AE raw};

\node[rep, minimum width=0.55cm, minimum height=0.36cm, text width=0.55cm] (leg2) at (5.0,-8.25) {};
\node[right=0.08cm of leg2, pub/legend] {AE representation};

\node[pp, minimum width=0.55cm, minimum height=0.36cm, text width=0.55cm] (leg3) at (10.0,-8.25) {};
\node[right=0.08cm of leg3, pub/legend] {Process / physics};

\node[config, minimum width=0.55cm, minimum height=0.36cm, text width=0.55cm] (leg4) at (1.0,-9.0) {};
\node[right=0.08cm of leg4, pub/legend] {Config. category};

\node[model, minimum width=0.55cm, minimum height=0.36cm, text width=0.55cm] (leg5) at (5.0,-9.0) {};
\node[right=0.08cm of leg5, pub/legend] {Model};

\node[target, minimum width=0.55cm, minimum height=0.36cm, text width=0.55cm] (leg6) at (10.0,-9.0) {};
\node[right=0.08cm of leg6, pub/legend] {Target};

\end{tikzpicture}%
}

\caption{End-to-end experimental workflow. Raw AE and vibration signals
are converted into physics-aware representations, optionally combined
with process parameters and physics-informed features, and evaluated by
auditable classical or deep-learning models under the unified LOGO protocol. The
two AE spectrogram planes map to the physical narrowband and broadband raw
AE columns; their analogue filter specifications and gains were not
archived. The full 27-model, 38-input inventory is given in the
Supplementary Materials.}
\label{fig:workflow}
\end{figure}

Software versions used for the experiments are Python 3.11, NumPy 1.26,
SciPy 1.14, scikit-learn 1.5, LightGBM 4.5, XGBoost 2.1, PyTorch 2.4,
Kymatio 0.3.0, and SHAP 0.51.0.

\subsection{Configuration abbreviations}\label{sec:config-abbreviations}

Table~\ref{tbl:config-abbreviations} lists the short labels used for
input configurations in the Results tables and figures. The labels follow
the convention \texttt{modality--representation} (e.g.~AE-dB-z for the
acoustic-emission per-sample-standardised dB spectrogram), with fused inputs joined by a plus
sign and auxiliary feature groups (physics-informed features, process
parameters) indicated by \texttt{phys} and \texttt{PP}.

\begin{table}[htbp]
\centering
\caption{Abbreviations for model input configurations used in the
Results section.}\label{tbl:config-abbreviations}
\begin{tabular}[]{@{}ll@{}}
\toprule
Abbreviation & Full configuration \\
\midrule
AE-dB-z & AE per-sample-standardised dB spectrogram \\
Vib-dB-z & Vibration per-sample-standardised dB spectrogram \\
AE-dB-z + Vib-dB-z & Pass-level fused AE and vibration dB-z descriptors \\
AE-dB-z + Vib-dB-z + PP & Pass-level fused dB-z descriptors + process parameters \\
AE-log-mel & AE log-mel spectrogram \\
Vib-log-mel & Vibration log-mel spectrogram \\
AE-log-mel + Vib-log-mel & Pass-level fused AE and vibration log-mel descriptors \\
AE-log-mel + Vib-log-mel + PP & Pass-level fused log-mel descriptors + process parameters \\
AE-dB & AE dB spectrogram \\
Vib-dB & Vibration dB spectrogram \\
AE-dB + Vib-dB & Pass-level fused AE and vibration dB descriptors \\
AE-dB + Vib-dB + PP & Pass-level fused dB descriptors + process parameters \\
AE-dB + Vib-dB + phys & Pass-level fused dB descriptors + physics-informed features \\
AE-dB + Vib-dB + phys + PP & Pass-level fused dB descriptors + physics + process parameters \\
Vib-WST & Vibration wavelet scattering transform \\
Vib-traj & Vibration trajectory/image representation \\
AE/Vib feats + phys + PP & AE and vibration hand-crafted features + physics + process parameters \\
Process params & Process parameters only \\
\bottomrule
\end{tabular}
\end{table}

In internal file names and cached feature files the full-resolution
dB-z descriptors are labelled \texttt{ae\_logspec} and
\texttt{vib\_logspec}, which correspond one-to-one to the AE-dB-z and
Vib-dB-z labels used here and in the result tables. The 64-bin log-mel
representations are stored as \texttt{ae\_mel} and \texttt{vib\_mel}.
